\documentclass[12pt,a4paper]{report}

\usepackage[margin=2.5cm]{geometry}
\usepackage[T1]{fontenc}
\usepackage[utf8]{inputenc}
\usepackage{lmodern}
\usepackage{graphicx}
\usepackage{float}
\usepackage{longtable}
\usepackage{array}
\usepackage{booktabs}
\usepackage{xcolor}
\usepackage{xurl}
\usepackage{hyperref}
\usepackage{listings}
\usepackage{caption}
\usepackage{enumitem}
\usepackage{titlesec}

\hypersetup{
  colorlinks=true,
  linkcolor=black,
  citecolor=black,
  urlcolor=blue,
  pdfauthor={To be completed},
  pdftitle={Prompt\_Env Technical Project Report}
}

\definecolor{codebg}{RGB}{248,248,248}
\definecolor{accent}{RGB}{35,110,55}
\definecolor{border}{RGB}{210,210,210}

\lstset{
  basicstyle=\ttfamily\small,
  backgroundcolor=\color{codebg},
  frame=single,
  rulecolor=\color{border},
  breaklines=true,
  columns=fullflexible,
  keepspaces=true,
  showstringspaces=false,
  tabsize=2
}

\setlist[itemize]{leftmargin=1.2cm}
\setlist[enumerate]{leftmargin=1.2cm}
\renewcommand{\arraystretch}{1.18}

\newcommand{\ProjectName}{Prompt\_Env}
\newcommand{\StudentName}{To be completed}
\newcommand{\StudentId}{To be completed}
\newcommand{\SupervisorName}{To be completed}
\newcommand{\GitHubUrl}{To be completed}
\newcommand{\RailwayUrl}{https://prompt-engineering-platform-production.up.railway.app}

\newcommand{\ReportFigure}[4]{
  \begin{figure}[H]
    \centering
    \IfFileExists{#1}{
      \includegraphics[width=0.95\textwidth]{#1}
    }{
      \fbox{\parbox{0.88\textwidth}{
        \centering
        \vspace{1.2cm}
        Missing image: \path{#1}\\
        #4
        \vspace{1.2cm}
      }}
    }
    \caption{#2}
    \label{#3}
  \end{figure}
}

\begin{document}

\begin{titlepage}
  \centering
  \vspace*{1cm}
  {\Large Institut National des Postes et Telecommunications (INPT)\par}
  \vspace{0.4cm}
  {\large Advanced Software Engineering and Data Science (ASEDS)\par}
  \vspace{1.6cm}
  {\Huge\bfseries Technical Project Report\par}
  \vspace{0.7cm}
  {\LARGE\bfseries \ProjectName: Prompt Engineering Platform\par}
  \vspace{1.5cm}
  \begin{tabular}{>{\bfseries}p{5cm}p{8cm}}
    Student name(s): & \StudentName \\
    Student ID(s): & \StudentId \\
    Supervisor / Encadrant: & \SupervisorName \\
    Academic year: & 2025/2026 \\
    Submission date: & \today \\
    GitHub URL: & \GitHubUrl \\
    Railway URL: & \url{\RailwayUrl} \\
  \end{tabular}
  \vfill
  {\large Full-stack prompt engineering, evaluation, and experiment tracking platform\par}
\end{titlepage}

\pagenumbering{roman}
\chapter*{Abstract}
\addcontentsline{toc}{chapter}{Abstract}
\ProjectName{} is a full-stack prompt engineering platform that allows users to register large language model providers, design prompts, version prompt templates, run inference, evaluate outputs, and compare experiments across models and datasets. The project addresses the gap between ad hoc prompt testing and professional prompt lifecycle management by treating prompts as versioned, testable, and measurable software artifacts.

The implemented system combines a React and Vite single-page application, a FastAPI backend, SQLAlchemy models, Alembic migrations, a Supabase PostgreSQL database, and a Railway deployment target. It supports multiple artificial intelligence providers, including Anthropic, OpenAI, Google Gemini, Groq, Mistral, HuggingFace, and custom OpenAI-compatible endpoints. API keys are encrypted at rest using Fernet and are only decrypted server-side when an inference request is executed.

The main outcomes are a secure backend inference layer, a prompt library with per-prompt version histories, dataset-driven batch evaluation, AI-assisted scoring with custom metrics, experiment logging, and comparison views for batch runs. The project also produced a realistic set of engineering lessons: network constraints around managed databases, deployment cold starts, prompt editor cursor synchronization, rate limit handling, JSON parsing reliability for evaluator models, and the importance of modular frontend structure.

\tableofcontents
\listoffigures
\listoftables
\clearpage
\pagenumbering{arabic}

\chapter{Introduction}

\section{Context}
Large language models (LLMs) are increasingly used in customer support, search, summarization, coding, data analysis, and agentic workflows. In these systems, the prompt is no longer a casual text instruction. It becomes a reusable interface between the application and the model. Small changes to system prompts, user templates, variables, or model parameters can change output quality, latency, and cost.

Modern teams therefore need prompt engineering platforms that provide versioning, repeatable tests, datasets, evaluation metrics, and experiment history. Commercial and open-source platforms such as LangSmith, Braintrust, PromptLayer, and Langfuse show that prompt iteration is now part of the software engineering workflow, not a separate manual activity.

\section{Problem Statement}
The initial workflow for testing prompts is usually fragmented: prompts are copied between notebooks, provider dashboards, local files, and chat interfaces. This makes it difficult to answer basic engineering questions:
\begin{itemize}
  \item Which prompt version produced this output?
  \item Which model and parameters were used?
  \item Did a new version improve relevance or correctness?
  \item How does a prompt behave across a full dataset instead of one example?
  \item Are API keys protected from browser exposure?
\end{itemize}

\ProjectName{} fills this gap by providing a single application for prompt creation, version control, inference, evaluation, and experiment comparison.

\section{Objectives}
The project objectives were:
\begin{enumerate}
  \item Build a usable prompt engineering workspace for prompt creation, editing, and versioning.
  \item Support at least five AI providers through one normalized backend inference layer.
  \item Store provider API keys securely and prevent client-side exposure.
  \item Allow prompts to be tested against datasets with variable interpolation.
  \item Log each run as an experiment with prompt, model, dataset, latency, token, cost, and score metadata.
  \item Provide AI-assisted scoring using custom workspace-level metrics.
  \item Deploy the backend to a cloud environment and connect it to a managed PostgreSQL database.
\end{enumerate}

\section{Scope}
The implemented scope includes authentication, workspace creation, model registry, prompt library, prompt studio, datasets, experiments, evaluations, custom metrics, batch runs, batch scoring, and workspace settings. The project does not yet include multi-user invitation workflows, a deployed public frontend, refresh tokens, persistent background-job storage, or complete DevOps automation.

\section{Report Structure}
This report presents the state of the art, requirements, architecture, database design, backend implementation, frontend implementation, evaluation system, security, testing approach, challenges, results, future work, conclusion, references, and appendices.

\chapter{State of the Art and Literature Review}

\section{Existing Platforms}
LangSmith is a framework-agnostic platform for building, debugging, evaluating, and deploying LLM applications. It includes observability, evaluation, prompt engineering, and deployment workflows. Braintrust focuses on AI observability, evaluation, prompt iteration, production logs, datasets, experiments, and continuous improvement. PromptLayer provides prompt registry, evaluations, dataset management, observability, release labels, A/B testing, and analytics. Langfuse is an open-source AI engineering platform for observability, prompt management, evaluation, metrics, and self-hosting.

\begin{longtable}{p{0.18\textwidth}p{0.15\textwidth}p{0.15\textwidth}p{0.15\textwidth}p{0.15\textwidth}p{0.15\textwidth}}
\caption{Comparison of prompt engineering and LLM evaluation platforms.}\\
\toprule
Feature & LangSmith & Braintrust & PromptLayer & Langfuse & \ProjectName{} \\
\midrule
\endfirsthead
\toprule
Feature & LangSmith & Braintrust & PromptLayer & Langfuse & \ProjectName{} \\
\midrule
\endhead
Prompt versioning & Yes & Yes & Yes & Yes & Yes \\
Dataset-driven tests & Yes & Yes & Yes & Yes & Yes \\
Experiment tracking & Yes & Yes & Yes & Yes & Yes \\
LLM-as-judge scoring & Yes & Yes & Yes & Yes & Yes \\
Custom metrics & Yes & Yes & Yes & Yes & Yes \\
Multi-provider registry & Yes & Yes & Yes & Yes & Yes \\
Encrypted user API keys & Platform-dependent & Platform-dependent & Platform-dependent & Platform-dependent & Implemented with Fernet \\
Open-source/self-hostable & No, managed plus options & Enterprise options & Self-hosted option & Yes & Codebase-owned project \\
Frontend deployment & Yes & Yes & Yes & Yes & Not yet \\
Local educational implementation & No & No & No & No & Yes \\
\bottomrule
\end{longtable}

\section{Key Concepts}
Prompt engineering is the process of designing model instructions, examples, variables, and output constraints to obtain reliable LLM behavior. LLM evaluation measures the quality, safety, and task fit of outputs, using human review, rule-based checks, model-based scoring, or combinations of these methods. Versioning records prompt changes over time and makes experiments reproducible. Dataset-driven testing evaluates prompts over many input examples rather than isolated manual trials.

\section{Project Differentiation}
\ProjectName{} is intentionally smaller than commercial products, but it is educationally valuable because the full architecture is visible. The project implements the core lifecycle from scratch: authentication, model registry, encrypted keys, prompt versioning, dataset storage, inference routing, scoring, and comparison. This gives the project a transparent engineering narrative that is well suited to a technical report and demonstration.

\chapter{Requirements Analysis}

\section{Functional Requirements}
\begin{longtable}{p{0.15\textwidth}p{0.78\textwidth}}
\caption{Functional requirements.}\\
\toprule
ID & Requirement \\
\midrule
\endfirsthead
\toprule
ID & Requirement \\
\midrule
\endhead
FR-01 & Users must be able to register, log in, and retrieve their authenticated profile. \\
FR-02 & A workspace must be created automatically during registration. \\
FR-03 & Users must be able to register, update, validate, activate, deactivate, and delete AI models. \\
FR-04 & The system must support multiple providers: Anthropic, OpenAI, Google, Groq, Mistral, HuggingFace, and custom endpoints. \\
FR-05 & Users must be able to create prompts with name, description, tags, system prompt, and user template. \\
FR-06 & Prompts must support version history with per-prompt version numbers and commit messages. \\
FR-07 & Users must be able to save a version in place and commit a new version explicitly. \\
FR-08 & User templates must support variable interpolation with tokens such as \texttt{\{customer\_issue\}}. \\
FR-09 & Users must be able to create, import, edit, export, and delete datasets. \\
FR-10 & The system must run single prompt inference and store each result as an experiment. \\
FR-11 & The system must run batch evaluation against dataset rows. \\
FR-12 & Users must be able to filter, sort, inspect, export, update, and delete experiments. \\
FR-13 & Users must be able to define custom evaluation metrics. \\
FR-14 & Users must be able to score one experiment or a batch of experiments using an AI scorer model. \\
FR-15 & Users must be able to compare named batch runs row by row. \\
FR-16 & Users must be able to update workspace and profile settings. \\
\bottomrule
\end{longtable}

\section{Non-Functional Requirements}
\begin{longtable}{p{0.15\textwidth}p{0.78\textwidth}}
\caption{Non-functional requirements.}\\
\toprule
ID & Requirement \\
\midrule
\endfirsthead
\toprule
ID & Requirement \\
\midrule
\endhead
NFR-01 & Provider API keys must never be stored in plaintext. \\
NFR-02 & Provider API keys must never be returned to the browser after initial submission. \\
NFR-03 & Backend endpoints must require JSON Web Token (JWT) authentication except registration and login. \\
NFR-04 & The application must support at least five AI providers. \\
NFR-05 & External HTTP calls must protect against Server-Side Request Forgery (SSRF). \\
NFR-06 & Authentication routes must include rate limiting. \\
NFR-07 & Prompt editing should preserve draft state across page reloads. \\
NFR-08 & Dataset cell edits should not block unrelated cells from being saved. \\
NFR-09 & Batch scoring must continue server-side if the user navigates away. \\
NFR-10 & The frontend should remain usable during Railway backend cold starts. \\
\bottomrule
\end{longtable}

\section{Use Case Diagram}
\ReportFigure{figures/diagrams/diagram_04_use_case.png}{Use case diagram for authenticated platform workflows.}{fig:use-case}{Render Diagram 4 from \path{project_diagrams.md}.}

\section{Use Case Descriptions}
\begin{longtable}{p{0.17\textwidth}p{0.20\textwidth}p{0.26\textwidth}p{0.26\textwidth}}
\caption{Major use case descriptions.}\\
\toprule
Use case & Preconditions & Main flow & Alternative flows \\
\midrule
\endfirsthead
\toprule
Use case & Preconditions & Main flow & Alternative flows \\
\midrule
\endhead
Authenticate & User has account or wants to register & Submit credentials, backend validates or creates account, JWT is returned & Invalid password, duplicate email, rate limit exceeded \\
Register model & User is authenticated & Enter provider, model ID, endpoint, key, and parameters; backend encrypts key & Endpoint blocked by SSRF policy, key validation fails \\
Create prompt & User is authenticated & Create prompt metadata, backend creates prompt and initial v1 & Missing name, backend validation error \\
Run prompt & Prompt version and active model exist & Interpolate variables, call provider, log experiment & Provider failure logs error experiment \\
Import dataset & User is authenticated & Upload CSV or JSON, preview, create rows & Invalid file format or empty rows \\
Score experiment & Experiment and scorer model exist & Build scoring prompt, call provider, parse JSON, update scores & Rate limit retry, parse retry, missing expected output \\
Compare batches & Two batch runs exist & Match experiments by row index and show side-by-side outputs & Missing corresponding row in one batch \\
\bottomrule
\end{longtable}

\chapter{System Architecture}

\section{High-Level Architecture}
\ReportFigure{figures/diagrams/diagram_01_system_architecture.png}{High-level architecture of the React frontend, FastAPI backend, Supabase database, and AI providers.}{fig:system-architecture}{Render Diagram 1 from \path{project_diagrams.md}.}

The application follows a three-tier architecture. The frontend is a React single-page application. The backend is a FastAPI service exposing REST endpoints under \path{/api}. The database is hosted PostgreSQL on Supabase. External AI providers are called only from the backend.

\section{Deployment Architecture}
\ReportFigure{figures/diagrams/diagram_02_deployment.png}{Deployment architecture from local development to Railway and Supabase.}{fig:deployment}{Render Diagram 2 from \path{project_diagrams.md}.}

Railway deploys the backend from the \path{backend/} directory using the Dockerfile. Supabase hosts the PostgreSQL database. The frontend currently runs locally through Vite on port 5173 and targets the Railway backend through \texttt{VITE\_API\_BASE\_URL} or the production URL fallback.

\section{Data Flow}
\ReportFigure{figures/diagrams/diagram_14_data_flow_dfd.png}{Level 1 data flow diagram.}{fig:data-flow}{Render Diagram 14 from \path{project_diagrams.md}.}

\section{Architectural Decisions}
\begin{longtable}{p{0.24\textwidth}p{0.68\textwidth}}
\caption{Architectural decisions and rationale.}\\
\toprule
Decision & Rationale \\
\midrule
\endfirsthead
\toprule
Decision & Rationale \\
\midrule
\endhead
FastAPI over Django or Flask & FastAPI provides automatic OpenAPI documentation, Pydantic validation, dependency injection, async HTTP support, and a smaller structure than Django for an API-first project. \\
Supabase over self-hosted PostgreSQL & Supabase provides managed PostgreSQL, dashboard visibility, and a free tier. The trade-off is network dependency, especially when port 5432 is blocked. \\
Railway for deployment & Railway supports simple GitHub-linked deployment, Dockerfile-based backend deployment, and environment variable management. \\
Backend-only AI calls & Moving provider calls to the backend prevents browser exposure of API keys and centralizes provider-specific request handling. \\
JWT over server sessions & JWTs simplify stateless API authentication for the deployed backend and local frontend. \\
Drafts separate from versions & Drafts are ephemeral and local to a user/device, while committed prompt versions are durable, shared application data. \\
Per-prompt version numbers & Version numbers are meaningful inside a prompt history and avoid misleading global version counters. \\
\bottomrule
\end{longtable}

\section{Backend Component Diagram}
\ReportFigure{figures/diagrams/diagram_13_backend_components.png}{Backend component diagram showing routers, services, and models.}{fig:backend-components}{Render Diagram 13 from \path{project_diagrams.md}.}

\section{Frontend Page Navigation}
\ReportFigure{figures/diagrams/diagram_12_frontend_navigation.png}{Frontend page navigation and authenticated app shell.}{fig:frontend-navigation}{Render Diagram 12 from \path{project_diagrams.md}.}

\chapter{Database Design}

\section{Entity Relationship Diagram}
\ReportFigure{figures/diagrams/diagram_03_entity_relationship.png}{Entity relationship diagram for the ten core database tables.}{fig:erd}{Render Diagram 3 from \path{project_diagrams.md}.}

\section{Table Descriptions}
\begin{longtable}{p{0.22\textwidth}p{0.34\textwidth}p{0.34\textwidth}}
\caption{Database table descriptions.}\\
\toprule
Table & Purpose & Key design details \\
\midrule
\endfirsthead
\toprule
Table & Purpose & Key design details \\
\midrule
\endhead
\texttt{users} & Stores registered users & Unique email, bcrypt password hash, display name, role \\
\texttt{workspaces} & Groups user-owned resources & Owner foreign key, created automatically at registration \\
\texttt{workspace\_members} & Future collaboration membership & Unique workspace-user pair, role field \\
\texttt{models} & Stores provider model configuration & Encrypted API key, provider, endpoint, model ID, parameters \\
\texttt{prompts} & Stores prompt metadata & Name, description, tags, creator, workspace ownership \\
\texttt{prompt\_versions} & Stores durable prompt content & Per-prompt version number, system prompt, user template, commit message \\
\texttt{datasets} & Stores dataset metadata & Columns stored as text array, category, version, workspace ownership \\
\texttt{dataset\_rows} & Stores dataset row payloads & Separate JSONB rows allow efficient list views with row counts \\
\texttt{experiments} & Stores run history and evaluation data & Denormalized snapshots preserve reproducibility after later edits \\
\texttt{evaluation\_metrics} & Stores workspace scoring rubrics & Supports inverse metrics such as Toxicity \\
\bottomrule
\end{longtable}

\section{Design Decisions}
Rows are stored separately from dataset metadata because dataset lists only need row counts and column previews, while detail views need full row data. Experiments store denormalized prompt, model, provider, variable, output, cost, and score snapshots so that historical runs remain understandable even after prompts or models are edited. Version numbers are scoped by prompt and protected by the unique constraint \texttt{uq\_prompt\_versions\_prompt\_version}.

\section{Migration Strategy}
Schema changes are managed with Alembic. The initial migration creates the main tables and constraints. A later migration adds \texttt{batch\_id}, \texttt{batch\_name}, and an index on batch identifiers for grouped experiment comparison.

\chapter{Backend Implementation}

\section{Technology Choices}
The backend uses FastAPI for API routes, SQLAlchemy for object-relational mapping (ORM), Alembic for migrations, \texttt{httpx} for asynchronous provider calls, \texttt{python-jose} for JWT handling, \texttt{passlib[bcrypt]} for password hashing, \texttt{cryptography} for Fernet encryption, \texttt{slowapi} for rate limiting, and Pydantic for validation and camelCase response models.

\section{Project Structure}
\begin{lstlisting}
backend/
  app/
    main.py              # FastAPI app, CORS, rate limit handler, router mounting
    database.py          # SQLAlchemy engine, session, Base
    core/
      auth.py            # Password hashing, JWT creation, user dependency
      config.py          # Environment-backed settings
      security.py        # SSRF endpoint validation
    models/              # SQLAlchemy database models
    schemas/             # Pydantic request/response schemas
    routers/             # REST endpoint groups
    services/
      ai_router.py       # Provider dispatcher
      crypto.py          # Fernet encrypt/decrypt helpers
      cost.py            # Token cost estimation
      retry.py           # Exponential backoff helpers
      experiment_logger.py
  alembic/               # Migrations
\end{lstlisting}

\section{API Design}
The backend exposes REST-style routes under \path{/api}. The frontend sends and receives camelCase JSON, while backend models use snake_case. Pydantic alias generators bridge these styles. Protected endpoints depend on \texttt{get\_current\_user}, which validates the Bearer token, decodes the JWT, and loads the user from the database.

\section{API Endpoint Map}
\ReportFigure{figures/diagrams/diagram_10_api_endpoint_map.png}{API endpoint map by FastAPI router.}{fig:endpoint-map}{Render Diagram 10 from \path{project_diagrams.md}.}

\section{Authentication System}
\ReportFigure{figures/diagrams/diagram_06_auth_sequence.png}{Registration, login, and protected request sequence.}{fig:auth-sequence}{Render Diagram 6 from \path{project_diagrams.md}.}

Registration creates the user, hashes the password, creates a workspace, adds an admin workspace membership, seeds default evaluation metrics, and returns a JWT. Login verifies the password and returns a JWT plus user and workspace data. The current implementation rate-limits registration to \texttt{5/hour}, login to \texttt{10/minute}, and profile fetches to \texttt{60/minute}.

\section{Multi-Provider Inference}
\ReportFigure{figures/diagrams/diagram_05_inference_sequence.png}{Inference sequence from browser action to provider call and experiment logging.}{fig:inference-sequence}{Render Diagram 5 from \path{project_diagrams.md}.}

The service \path{backend/app/services/ai_router.py} exposes a single \texttt{call\_provider(model, system\_prompt, user\_message)} function. It decrypts the model API key, validates the endpoint, dispatches to a provider-specific implementation, and returns normalized output, token, and latency metadata. OpenAI-compatible providers share a common handler. Google Gemini uses a separate body format and key placement. Anthropic uses the messages endpoint and \texttt{x-api-key} header.

\section{Retry System}
\texttt{retry.py} defines \texttt{RateLimitError}, \texttt{ParseError}, and \texttt{with\_exponential\_backoff}. Rate-limit retries use larger delays and parse retries use shorter delays because malformed model output usually resolves with another request.

\section{API Key Encryption}
Provider API keys are encrypted with Fernet before storage in \texttt{models.api\_key\_encrypted}. The Fernet key is supplied through \texttt{ENCRYPTION\_KEY}. API responses return masked values rather than plaintext keys.

\section{Security Architecture}
\ReportFigure{figures/diagrams/diagram_11_security_architecture.png}{Security architecture for encrypted API keys, JWTs, CORS, and provider calls.}{fig:security-architecture-backend}{Render Diagram 11 from \path{project_diagrams.md}.}

\section{Cost Estimation}
\texttt{cost.py} stores per-token rates for known models and returns \texttt{0.0} for unknown or free-tier models. Experiments persist cost as a floating-point value, avoiding previous crashes caused by string-formatted costs.

\chapter{Frontend Implementation}

\section{Technology Choices}
The frontend uses React 18, Vite, custom CSS and Tailwind utility configuration, IBM Plex typography, \texttt{lucide-react} icons, and browser \texttt{localStorage} for session and draft state. No external component library is used, which keeps the visual system custom and lightweight.

\section{Project Structure}
\begin{lstlisting}
frontend/src/
  App.jsx                       # Auth state, app shell, view routing
  components/
    Sidebar.jsx
    TopBar.jsx
  pages/
    LandingPage.jsx
    AuthPage.jsx
    ModelsPage.jsx
    PromptsPage.jsx
    PromptStudio.jsx
    DatasetsPage.jsx
    ExperimentsPage.jsx
    EvaluationsPage.jsx
    WorkspaceSettingsPage.jsx
  utils/
    api.js                      # Backend client
    auth.js                     # Token/user/workspace localStorage helpers
    helpers.js                  # Variable extraction, formatting helpers
    constants.js
\end{lstlisting}

\section{Design System}
\begin{longtable}{p{0.23\textwidth}p{0.20\textwidth}p{0.43\textwidth}}
\caption{Visual design tokens.}\\
\toprule
Token & Value & Usage \\
\midrule
\endfirsthead
\toprule
Token & Value & Usage \\
\midrule
\endhead
\texttt{--background} & \texttt{\#0f0e0d} & Main page background \\
\texttt{--surface} & \texttt{\#161613} & Cards and panels \\
\texttt{--surface-raised} & \texttt{\#1a1916} & Elevated elements \\
\texttt{--border} & \texttt{\#252320} & Borders and dividers \\
\texttt{--text-primary} & \texttt{\#f0ece4} & Main text \\
\texttt{--text-muted} & \texttt{\#6b6860} & Secondary text \\
\texttt{--accent} & \texttt{\#88d273} & Primary green action and status color \\
\texttt{--error} & \texttt{\#ff6b6b} & Error states \\
\texttt{--warning} & \texttt{\#e8a847} & Warning and high-latency indication \\
\bottomrule
\end{longtable}

The UI uses a warm dark theme, compact cards, small uppercase mono labels, version pills, stat chips, skeleton loaders, modals, drawers, and empty states. Latency is colored amber when slow but not red, because red is reserved for actual errors.

\section{Key Implementation Details}
\textbf{Variable highlighting.} The Prompt Studio user template uses a textarea and a mirrored highlight layer. Variable spans only change background and color, preventing cursor desynchronization. Scroll position is synchronized between textarea and mirror.

\textbf{Draft persistence.} Prompt drafts are stored in \texttt{pe\_drafts} in \texttt{localStorage}, keyed by prompt ID. Autosave uses an 800 ms debounce.

\textbf{Dataset editing.} Dataset cells use independent edit buffers and timers, so editing one cell does not block another. Pending edits are flushed before row and column operations.

\textbf{Modal behavior.} Modals close on Escape and outside click. Outside-click detection uses mouse down and mouse up tracking to avoid accidental close when dragging from inside to outside.

\textbf{Authentication flow.} Tokens are stored in \texttt{pe\_auth\_token}. The frontend decodes JWT expiry client-side for initial auth state and clears session on 401 responses.

\section{Landing Page}
The landing page explains the project value proposition, shows provider support, and routes unauthenticated users to login/register while authenticated users can return to the app. Fade-in animation is implemented with \texttt{IntersectionObserver}.

\chapter{AI Evaluation System}

\section{Overview}
The evaluation system converts raw prompt runs into measurable experiments. It supports manual scoring, AI scoring, custom metrics, expected-output injection from datasets, batch scoring jobs, and batch comparison.

\section{Prompt Lifecycle}
\ReportFigure{figures/diagrams/diagram_08_prompt_lifecycle.png}{Prompt lifecycle from creation to evaluation.}{fig:prompt-lifecycle}{Render Diagram 8 from \path{project_diagrams.md}.}

\section{Experiment State}
\ReportFigure{figures/diagrams/diagram_09_experiment_state.png}{Experiment state transitions during inference and scoring.}{fig:experiment-state}{Render Diagram 9 from \path{project_diagrams.md}.}

\section{Scoring System Evolution}
\begin{longtable}{p{0.16\textwidth}p{0.28\textwidth}p{0.42\textwidth}}
\caption{Evolution of the AI scoring system.}\\
\toprule
Phase & Problem addressed & Solution \\
\midrule
\endfirsthead
\toprule
Phase & Problem addressed & Solution \\
\midrule
\endhead
Phase 1 & Scores lacked ground truth context & Expected output is fetched from the dataset row when available \\
Phase 2 & Zero-shot JSON scoring was unreliable & Chain-of-thought reasoning is requested before final JSON \\
Phase 3 & Fixed metrics were too restrictive & Workspace custom metrics are stored and injected dynamically \\
Phase 4 & Frontend scoring loop stopped on navigation & FastAPI background task handles batch scoring server-side \\
\bottomrule
\end{longtable}

\section{Batch Scoring Sequence}
\ReportFigure{figures/diagrams/diagram_07_batch_scoring_sequence.png}{Asynchronous batch AI scoring sequence with frontend polling.}{fig:batch-scoring}{Render Diagram 7 from \path{project_diagrams.md}.}

\section{Scoring Prompt Template}
\begin{lstlisting}
You are an expert evaluator of LLM outputs.

Input:
<user_input>{interpolated_prompt}</user_input>

Model output:
<model_output>{output}</model_output>

Expected output, if available:
<expected_output>{expected_output}</expected_output>

Evaluate the output using the following metrics:
<metric name="Relevance">...</metric>
<metric name="Correctness">...</metric>
<metric name="Fluency">...</metric>
<metric name="Toxicity">...</metric>

Step 1: Reason briefly about each metric.
Step 2: Return final JSON only at the end:
{
  "scores": {"Metric": 0-10},
  "reasoning": {"Metric": "short explanation"}
}
\end{lstlisting}

Metric descriptions are wrapped in XML-style delimiters and sanitized to reduce prompt injection risk. The parser searches from the end of the model response because the final JSON is expected after reasoning.

\section{JSON Parsing Strategy}
The parser tries complete JSON blocks from last to first, falls back to parsing from the last opening brace, and finally uses regex-style extraction for individual metric scores. Full failure raises a parse error that can be retried.

\section{Custom Metrics}
Custom metrics are stored in \texttt{evaluation\_metrics}. Each metric has a name, description, default flag, order index, and \texttt{is\_inverse} flag. Inverse metrics such as Toxicity are displayed but excluded from overall score averages.

\section{Comparison View}
The comparison tab selects two named batch runs and matches experiments by \texttt{dataset\_row\_index}. This makes comparisons stable even when batches use different models or prompt versions.

\chapter{Security}

\section{Threat Model}
\ReportFigure{figures/diagrams/diagram_11_security_architecture.png}{Threat model and sensitive-data flow.}{fig:security-threat-model}{Render Diagram 11 from \path{project_diagrams.md}.}

\section{Implemented Security Measures}
\begin{longtable}{p{0.28\textwidth}p{0.62\textwidth}}
\caption{Implemented security controls.}\\
\toprule
Control & Implementation \\
\midrule
\endfirsthead
\toprule
Control & Implementation \\
\midrule
\endhead
API key encryption & Fernet encryption before database storage \\
Password hashing & bcrypt through \texttt{passlib}, with compatible bcrypt pin \\
JWT authentication & \texttt{python-jose}, HS256, seven-day expiry \\
API key masking & Model responses never return plaintext keys \\
SSRF protection & Provider host allowlist and private IP range blocking \\
CORS restriction & Explicit origins, methods, and headers \\
Rate limiting & \texttt{slowapi} limits on register, login, and profile fetches \\
IDOR mitigation & Workspace ownership checks on protected job operations \\
Prompt injection reduction & Metric input length limits, keyword blocklist, XML delimiters \\
\bottomrule
\end{longtable}

\section{Security Audit Results}
\begin{longtable}{p{0.12\textwidth}p{0.15\textwidth}p{0.38\textwidth}p{0.18\textwidth}}
\caption{Security audit findings.}\\
\toprule
ID & Severity & Finding & Status \\
\midrule
\endfirsthead
\toprule
ID & Severity & Finding & Status \\
\midrule
\endhead
SEC-001 & Critical & Ensure \texttt{.env} files are not committed and audit git history for secrets. & Pending verification \\
SEC-002 & Critical & Provider endpoint SSRF risk before URL validation. & Fixed \\
SEC-003 & High & Insecure direct object reference on scoring job cancellation. & Fixed \\
SEC-004 & High & Login endpoint brute-force and bcrypt CPU-exhaustion risk. & Fixed in current code with rate limiting \\
SEC-005 & Medium & Single static Fernet key means key exposure compromises stored provider keys. & Accepted risk; future KMS/Vault \\
SEC-006 & Medium & Overly broad CORS policy risk. & Fixed \\
\bottomrule
\end{longtable}

\section{Risk Management}
The remaining accepted risk is the static Fernet key. In this project phase, Railway environment variables provide a practical secret store. A production system should use managed key rotation through a cloud key management service or vault.

\chapter{Testing}

\section{Testing Approach}
Testing was primarily manual and scenario-based. FastAPI Swagger UI was used to validate backend endpoints and request/response shapes. The frontend was tested through realistic workflows: registration, adding models, creating prompts, importing datasets, running prompts, logging experiments, scoring outputs, and comparing batches.

\section{End-to-End Scenario}
The main demo scenario is a customer support assistant:
\begin{enumerate}
  \item Create a dataset with customer issues, tone requirements, and expected support answers.
  \item Register at least two models, such as a Groq-hosted model and a Google Gemini model.
  \item Create a prompt named ``Customer Support Assistant'' with variables for customer issue and customer tier.
  \item Run a single prompt test and inspect the logged experiment.
  \item Run a batch over the dataset.
  \item Score all outputs with an evaluator model.
  \item Compare two batches produced by different models or prompt versions.
\end{enumerate}

\section{Performance Observations}
Actual measured values should be captured during the final demo because they depend on provider, network, Railway cold-start state, and token usage. Table \ref{tab:demo-metrics} is prepared for insertion of final measurements.

\begin{longtable}{p{0.20\textwidth}p{0.18\textwidth}p{0.15\textwidth}p{0.15\textwidth}p{0.16\textwidth}}
\caption{Demo performance measurements to capture before submission.}
\label{tab:demo-metrics}\\
\toprule
Scenario & Model & Latency ms & Tokens & Cost estimate \\
\midrule
\endfirsthead
\toprule
Scenario & Model & Latency ms & Tokens & Cost estimate \\
\midrule
\endhead
Single support prompt & To be captured & To be captured & To be captured & To be captured \\
Batch row 1 & To be captured & To be captured & To be captured & To be captured \\
AI scoring request & To be captured & To be captured & To be captured & To be captured \\
\bottomrule
\end{longtable}

\section{Known Issues}
Known limitations include no deployed frontend, no refresh-token revocation mechanism, in-memory background job state, no invite UI for workspace collaboration, no local model support from Railway, and no persistent queue for long-running scoring jobs.

\chapter{Challenges and Solutions}

\begin{longtable}{p{0.24\textwidth}p{0.25\textwidth}p{0.34\textwidth}}
\caption{Major project challenges and solutions.}\\
\toprule
Challenge & Root cause & Solution \\
\midrule
\endfirsthead
\toprule
Challenge & Root cause & Solution \\
\midrule
\endhead
Supabase blocked on school network & Outbound port 5432 restrictions & Use mobile hotspot or Supabase SQL editor for migrations \\
Railway cold start and 499 errors & Free-tier backend sleeps after inactivity & Ping health endpoint before loading workspace data \\
Prompt editor cursor desync & Highlight spans changed text layout & Use textarea plus zero-cost mirrored highlight layer \\
Dataset cell edit blocking & Shared save state blocked independent edits & Use per-cell timers and flush pending edits before structural operations \\
Malformed scoring JSON & LLM responses include prose or invalid JSON & Ask for reasoning then final JSON and parse from response end \\
Provider rate limits & Free-tier token limits and repeated calls & Consolidate metrics into one scorer request and retry with backoff \\
Google endpoint bug & Literal placeholder was not interpolated & Replace placeholder with concrete model ID before request \\
\texttt{costEstimate} crash & Mixed string and numeric cost values & Store cost as float and format only in frontend \\
Hooks order black screen & React hook called after conditional return & Move all hooks before conditional rendering \\
Large frontend file edits & Single file exceeded practical edit context & Refactor into pages, components, and utilities \\
\bottomrule
\end{longtable}

\chapter{Results and Demo}

\section{Main Page Screenshots}
\ReportFigure{figures/screenshots/01_landing_page.png}{Landing page.}{fig:ss-landing}{Capture the unauthenticated landing page.}
\ReportFigure{figures/screenshots/02_login_page.png}{Login page.}{fig:ss-login}{Capture the login tab of the authentication page.}
\ReportFigure{figures/screenshots/03_models_page.png}{Models page.}{fig:ss-models}{Capture the model registry after at least one model is configured.}
\ReportFigure{figures/screenshots/04_prompts_library.png}{Prompts library.}{fig:ss-prompts}{Capture the prompt library with at least one prompt.}
\ReportFigure{figures/screenshots/05_prompt_studio.png}{Prompt Studio.}{fig:ss-studio}{Capture a prompt version with variables and an output preview.}
\ReportFigure{figures/screenshots/06_datasets_page.png}{Datasets page.}{fig:ss-datasets}{Capture dataset cards.}
\ReportFigure{figures/screenshots/07_experiments_page.png}{Experiments page.}{fig:ss-experiments}{Capture the experiment table with filters.}
\ReportFigure{figures/screenshots/08_evaluations_overview.png}{Evaluations overview tab.}{fig:ss-eval-overview}{Capture aggregate evaluation statistics.}
\ReportFigure{figures/screenshots/09_evaluations_comparison.png}{Evaluations comparison tab.}{fig:ss-eval-comparison}{Capture two selected batch runs.}
\ReportFigure{figures/screenshots/10_evaluations_batch_eval.png}{Batch evaluation tab.}{fig:ss-batch-eval}{Capture the batch evaluation workflow.}
\ReportFigure{figures/screenshots/11_workspace_settings.png}{Workspace settings page.}{fig:ss-workspace}{Capture profile, workspace, and metric settings.}

\section{Demo Walkthrough}
The demo should show the customer support assistant from prompt creation to batch comparison:
\begin{enumerate}
  \item Register or log into the application.
  \item Add two provider models and validate at least one key.
  \item Create a support dataset with expected outputs.
  \item Create a prompt with variables such as \texttt{\{customer\_issue\}} and \texttt{\{tone\}}.
  \item Run the prompt once and inspect the experiment log.
  \item Run a batch evaluation against the dataset.
  \item Score the batch with AI.
  \item Run a second batch with another model or prompt version.
  \item Compare both batches in the comparison tab.
\end{enumerate}

\section{Batch Comparison Result}
\ReportFigure{figures/screenshots/12_demo_batch_comparison_two_models.png}{Comparison of two batch runs using different models or prompt versions.}{fig:ss-demo-comparison}{Capture the comparison tab after selecting two named batches.}

\chapter{Future Work}

\begin{itemize}
  \item Workspace collaboration with invitations, roles, and workspace switching.
  \item Local model support through Docker Compose self-hosting or secure tunneling.
  \item DevOps pipeline using GitHub Actions, Kubernetes, and Argo CD.
  \item Per-row scoring from the Batch Eval table.
  \item Scoring prompt transparency so users can inspect the evaluator prompt.
  \item HuggingFace model browsing and discovery.
  \item Persistent background jobs using Redis Queue, Celery, or a database-backed job table.
  \item Refresh tokens and token revocation.
  \item Key rotation through a managed KMS or vault.
  \item Frontend deployment to Vercel, Netlify, or Railway static hosting.
\end{itemize}

\chapter{Conclusion}

\ProjectName{} achieved its main objectives: it provides a complete prompt lifecycle from model registration and prompt versioning to dataset-driven experiments, AI scoring, and batch comparison. The backend centralizes sensitive provider calls, encrypts API keys, applies authentication, validates provider endpoints, and persists reproducible experiment snapshots. The frontend provides a coherent workspace for practical prompt engineering workflows.

The project also demonstrates the real complexity of production-facing LLM tooling. Reliable prompt platforms must handle security, provider differences, rate limits, malformed model output, latency, cost, reproducibility, and user experience. The result is a strong full-stack foundation that can be extended into a collaborative and deployable prompt engineering product.

\chapter{References}

\begin{enumerate}
  \item FastAPI documentation. \url{https://fastapi.tiangolo.com/}
  \item SQLAlchemy documentation. \url{https://docs.sqlalchemy.org/}
  \item Alembic documentation. \url{https://alembic.sqlalchemy.org/}
  \item React documentation. \url{https://react.dev/}
  \item Vite documentation. \url{https://vitejs.dev/}
  \item LangSmith documentation. \url{https://docs.langchain.com/langsmith/home}
  \item Braintrust documentation. \url{https://www.braintrust.dev/docs}
  \item PromptLayer documentation. \url{https://docs.promptlayer.com/overview}
  \item Langfuse documentation. \url{https://langfuse.com/docs}
  \item OWASP Top Ten. \url{https://owasp.org/www-project-top-ten/}
  \item OWASP Server-Side Request Forgery Prevention Cheat Sheet. \url{https://cheatsheetseries.owasp.org/cheatsheets/Server_Side_Request_Forgery_Prevention_Cheat_Sheet.html}
  \item CWE-918: Server-Side Request Forgery. \url{https://cwe.mitre.org/data/definitions/918.html}
  \item Anthropic API documentation. \url{https://docs.anthropic.com/}
  \item OpenAI API documentation. \url{https://platform.openai.com/docs/}
  \item Google Gemini API documentation. \url{https://ai.google.dev/gemini-api/docs}
\end{enumerate}

\appendix

\chapter{Full API Endpoint Reference}

\begin{longtable}{p{0.13\textwidth}p{0.28\textwidth}p{0.48\textwidth}}
\caption{API endpoint reference.}\\
\toprule
Method & Path & Description \\
\midrule
\endfirsthead
\toprule
Method & Path & Description \\
\midrule
\endhead
POST & \path{/api/auth/register} & Create user, workspace, metrics, and token \\
POST & \path{/api/auth/login} & Authenticate user and return token \\
GET & \path{/api/auth/me} & Return current user and workspace \\
PATCH & \path{/api/auth/me} & Update current display name \\
GET & \path{/api/models} & List workspace models \\
POST & \path{/api/models} & Create encrypted model configuration \\
PATCH & \path{/api/models/{id}} & Update model configuration \\
DELETE & \path{/api/models/{id}} & Delete model \\
POST & \path{/api/models/validate} & Validate provider credentials and endpoint \\
GET & \path{/api/prompts} & List prompts with filters \\
POST & \path{/api/prompts} & Create prompt and initial version \\
PATCH & \path{/api/prompts/{id}} & Update prompt metadata \\
POST & \path{/api/prompts/{id}/duplicate} & Duplicate prompt and versions \\
DELETE & \path{/api/prompts/{id}} & Delete prompt \\
GET & \path{/api/prompts/{id}/versions} & List prompt versions \\
POST & \path{/api/prompts/{id}/versions} & Commit new version \\
GET & \path{/api/prompts/{id}/versions/{versionId}} & Get one version \\
PATCH & \path{/api/prompts/{id}/versions/{versionId}} & Save version in place \\
GET & \path{/api/datasets} & List datasets \\
POST & \path{/api/datasets} & Create dataset and rows \\
GET & \path{/api/datasets/{id}} & Get dataset detail \\
PUT & \path{/api/datasets/{id}} & Replace dataset fields and rows \\
DELETE & \path{/api/datasets/{id}} & Delete dataset \\
POST & \path{/api/datasets/import} & Import parsed CSV or JSON \\
GET & \path{/api/experiments} & List experiments with filters \\
POST & \path{/api/experiments} & Create manual experiment \\
PATCH & \path{/api/experiments/{id}} & Update notes, tags, score, scores, reasoning \\
DELETE & \path{/api/experiments/{id}} & Delete one experiment \\
POST & \path{/api/experiments/bulk-delete} & Delete multiple experiments \\
GET & \path{/api/experiments/batches} & List grouped batch runs \\
PATCH & \path{/api/experiments/batches/{batchId}} & Rename batch \\
POST & \path{/api/inference/run} & Run prompt and log experiment \\
POST & \path{/api/evaluations/score} & Score one experiment \\
POST & \path{/api/evaluations/batch-run} & Run prompt against dataset rows \\
POST & \path{/api/evaluations/score-batch} & Start background batch scoring \\
GET & \path{/api/evaluations/score-batch/{jobId}} & Poll scoring progress \\
POST & \path{/api/evaluations/score-batch/{jobId}/cancel} & Cancel scoring job \\
GET & \path{/api/metrics} & List evaluation metrics \\
POST & \path{/api/metrics} & Create custom metric \\
PATCH & \path{/api/metrics/{id}} & Update metric \\
DELETE & \path{/api/metrics/{id}} & Delete metric \\
PATCH & \path{/api/workspaces/{id}} & Rename workspace \\
\bottomrule
\end{longtable}

\chapter{Database Schema DDL}

\begin{lstlisting}[language=SQL]
CREATE TABLE users (
  id UUID PRIMARY KEY,
  email TEXT UNIQUE NOT NULL,
  password_hash TEXT NOT NULL,
  name TEXT NOT NULL,
  role VARCHAR DEFAULT 'member' NOT NULL,
  created_at TIMESTAMPTZ DEFAULT now() NOT NULL
);

CREATE TABLE workspaces (
  id UUID PRIMARY KEY,
  name TEXT NOT NULL,
  owner_id UUID NOT NULL REFERENCES users(id) ON DELETE CASCADE,
  created_at TIMESTAMPTZ DEFAULT now() NOT NULL
);

CREATE TABLE workspace_members (
  id UUID PRIMARY KEY,
  workspace_id UUID NOT NULL REFERENCES workspaces(id) ON DELETE CASCADE,
  user_id UUID NOT NULL REFERENCES users(id) ON DELETE CASCADE,
  role VARCHAR DEFAULT 'member' NOT NULL,
  UNIQUE(workspace_id, user_id)
);

CREATE TABLE models (
  id UUID PRIMARY KEY,
  workspace_id UUID NOT NULL REFERENCES workspaces(id) ON DELETE CASCADE,
  name TEXT NOT NULL,
  provider TEXT NOT NULL,
  model_id TEXT NOT NULL,
  endpoint TEXT NOT NULL,
  api_key_encrypted TEXT NOT NULL,
  temperature FLOAT DEFAULT 0.7 NOT NULL,
  max_tokens INTEGER DEFAULT 1024 NOT NULL,
  top_p FLOAT DEFAULT 1.0 NOT NULL,
  stop_sequences TEXT[] DEFAULT '{}' NOT NULL,
  status VARCHAR DEFAULT 'active' NOT NULL,
  created_at TIMESTAMPTZ DEFAULT now() NOT NULL,
  updated_at TIMESTAMPTZ DEFAULT now() NOT NULL
);

CREATE TABLE prompts (
  id UUID PRIMARY KEY,
  workspace_id UUID NOT NULL REFERENCES workspaces(id) ON DELETE CASCADE,
  name TEXT NOT NULL,
  description TEXT DEFAULT '' NOT NULL,
  tags TEXT[] DEFAULT '{}' NOT NULL,
  created_by UUID REFERENCES users(id),
  created_at TIMESTAMPTZ DEFAULT now() NOT NULL,
  updated_at TIMESTAMPTZ DEFAULT now() NOT NULL
);

CREATE TABLE prompt_versions (
  id UUID PRIMARY KEY,
  prompt_id UUID NOT NULL REFERENCES prompts(id) ON DELETE CASCADE,
  version_number INTEGER NOT NULL,
  system_prompt TEXT DEFAULT '' NOT NULL,
  user_template TEXT DEFAULT '' NOT NULL,
  commit_message TEXT DEFAULT '' NOT NULL,
  created_by UUID REFERENCES users(id),
  created_at TIMESTAMPTZ DEFAULT now() NOT NULL,
  UNIQUE(prompt_id, version_number)
);
\end{lstlisting}

\begin{lstlisting}[language=SQL]
CREATE TABLE datasets (
  id UUID PRIMARY KEY,
  workspace_id UUID NOT NULL REFERENCES workspaces(id) ON DELETE CASCADE,
  name TEXT NOT NULL,
  category TEXT DEFAULT 'Custom' NOT NULL,
  version TEXT DEFAULT 'v1' NOT NULL,
  columns TEXT[] NOT NULL,
  created_by UUID REFERENCES users(id),
  created_at TIMESTAMPTZ DEFAULT now() NOT NULL,
  updated_at TIMESTAMPTZ DEFAULT now() NOT NULL
);

CREATE TABLE dataset_rows (
  id UUID PRIMARY KEY,
  dataset_id UUID NOT NULL REFERENCES datasets(id) ON DELETE CASCADE,
  row_index INTEGER NOT NULL,
  row_data JSONB NOT NULL
);

CREATE TABLE experiments (
  id UUID PRIMARY KEY,
  workspace_id UUID NOT NULL REFERENCES workspaces(id) ON DELETE CASCADE,
  prompt_id UUID REFERENCES prompts(id),
  prompt_version_id UUID REFERENCES prompt_versions(id),
  model_id UUID REFERENCES models(id),
  dataset_id UUID REFERENCES datasets(id),
  dataset_row_index INTEGER,
  batch_id TEXT,
  batch_name TEXT,
  prompt_name TEXT,
  prompt_version TEXT,
  model_name TEXT,
  provider TEXT,
  system_prompt TEXT,
  user_template TEXT,
  variable_values JSONB DEFAULT '{}'::jsonb NOT NULL,
  interpolated_prompt TEXT,
  output TEXT,
  latency_ms INTEGER,
  input_tokens INTEGER,
  output_tokens INTEGER,
  total_tokens INTEGER,
  cost_estimate FLOAT DEFAULT 0 NOT NULL,
  status TEXT DEFAULT 'success' NOT NULL,
  error_message TEXT,
  score FLOAT,
  scores JSONB DEFAULT '{}'::jsonb NOT NULL,
  reasoning JSONB DEFAULT '{}'::jsonb NOT NULL,
  tags TEXT[] DEFAULT '{}' NOT NULL,
  notes TEXT DEFAULT '' NOT NULL,
  created_by UUID REFERENCES users(id),
  created_at TIMESTAMPTZ DEFAULT now() NOT NULL
);

CREATE INDEX ix_experiments_batch_id ON experiments(batch_id);
\end{lstlisting}

\begin{lstlisting}[language=SQL]
CREATE TABLE evaluation_metrics (
  id UUID PRIMARY KEY,
  workspace_id UUID NOT NULL REFERENCES workspaces(id) ON DELETE CASCADE,
  name TEXT NOT NULL,
  description TEXT NOT NULL,
  is_inverse BOOLEAN DEFAULT false NOT NULL,
  is_default BOOLEAN DEFAULT true NOT NULL,
  order_index INTEGER DEFAULT 0 NOT NULL,
  created_at TIMESTAMPTZ DEFAULT now() NOT NULL
);
\end{lstlisting}

\chapter{Environment Variables Reference}

\begin{longtable}{p{0.32\textwidth}p{0.56\textwidth}}
\caption{Required backend environment variables.}\\
\toprule
Variable & Purpose \\
\midrule
\endfirsthead
\toprule
Variable & Purpose \\
\midrule
\endhead
\texttt{DATABASE\_URL} & PostgreSQL connection string with SSL mode \\
\texttt{SUPABASE\_URL} & Supabase project URL \\
\texttt{SUPABASE\_ANON\_KEY} & Supabase anonymous key \\
\texttt{SUPABASE\_SERVICE\_KEY} & Supabase service key \\
\texttt{SECRET\_KEY} & JWT signing secret \\
\texttt{ALGORITHM} & JWT algorithm, defaults to HS256 \\
\texttt{ACCESS\_TOKEN\_EXPIRE\_MINUTES} & JWT expiry, currently seven days \\
\texttt{ENCRYPTION\_KEY} & Fernet key for API key encryption \\
\texttt{ENVIRONMENT} & Runtime environment label \\
\texttt{ALLOWED\_ORIGINS} & Comma-separated CORS origins \\
\texttt{VITE\_API\_BASE\_URL} & Frontend build-time backend base URL \\
\bottomrule
\end{longtable}

\chapter{All Diagrams}

\ReportFigure{figures/diagrams/diagram_01_system_architecture.png}{Diagram 1: System architecture.}{fig:appendix-diagram-01}{Render Diagram 1 from \path{project_diagrams.md}.}
\ReportFigure{figures/diagrams/diagram_02_deployment.png}{Diagram 2: Deployment architecture.}{fig:appendix-diagram-02}{Render Diagram 2 from \path{project_diagrams.md}.}
\ReportFigure{figures/diagrams/diagram_03_entity_relationship.png}{Diagram 3: Entity relationship diagram.}{fig:appendix-diagram-03}{Render Diagram 3 from \path{project_diagrams.md}.}
\ReportFigure{figures/diagrams/diagram_04_use_case.png}{Diagram 4: Use case diagram.}{fig:appendix-diagram-04}{Render Diagram 4 from \path{project_diagrams.md}.}
\ReportFigure{figures/diagrams/diagram_05_inference_sequence.png}{Diagram 5: Inference flow.}{fig:appendix-diagram-05}{Render Diagram 5 from \path{project_diagrams.md}.}
\ReportFigure{figures/diagrams/diagram_06_auth_sequence.png}{Diagram 6: Authentication flow.}{fig:appendix-diagram-06}{Render Diagram 6 from \path{project_diagrams.md}.}
\ReportFigure{figures/diagrams/diagram_07_batch_scoring_sequence.png}{Diagram 7: Batch AI scoring.}{fig:appendix-diagram-07}{Render Diagram 7 from \path{project_diagrams.md}.}
\ReportFigure{figures/diagrams/diagram_08_prompt_lifecycle.png}{Diagram 8: Prompt lifecycle.}{fig:appendix-diagram-08}{Render Diagram 8 from \path{project_diagrams.md}.}
\ReportFigure{figures/diagrams/diagram_09_experiment_state.png}{Diagram 9: Experiment state.}{fig:appendix-diagram-09}{Render Diagram 9 from \path{project_diagrams.md}.}
\ReportFigure{figures/diagrams/diagram_10_api_endpoint_map.png}{Diagram 10: API endpoint map.}{fig:appendix-diagram-10}{Render Diagram 10 from \path{project_diagrams.md}.}
\ReportFigure{figures/diagrams/diagram_11_security_architecture.png}{Diagram 11: Security architecture.}{fig:appendix-diagram-11}{Render Diagram 11 from \path{project_diagrams.md}.}
\ReportFigure{figures/diagrams/diagram_12_frontend_navigation.png}{Diagram 12: Frontend navigation.}{fig:appendix-diagram-12}{Render Diagram 12 from \path{project_diagrams.md}.}
\ReportFigure{figures/diagrams/diagram_13_backend_components.png}{Diagram 13: Backend components.}{fig:appendix-diagram-13}{Render Diagram 13 from \path{project_diagrams.md}.}
\ReportFigure{figures/diagrams/diagram_14_data_flow_dfd.png}{Diagram 14: Level 1 data flow.}{fig:appendix-diagram-14}{Render Diagram 14 from \path{project_diagrams.md}.}

\chapter{Security Audit Full Report}
The security audit focused on secret handling, authentication, authorization, outbound HTTP calls, CORS policy, and prompt-injection vectors. The most important lesson is that LLM applications combine normal web security risks with provider-specific risks. A prompt engineering platform must defend both the user data plane and the model call plane.

\begin{itemize}
  \item \textbf{SEC-001}: Audit git history for accidental environment variable commits before submission or deployment handoff.
  \item \textbf{SEC-002}: SSRF was mitigated through provider host allowlists and private/reserved IP blocking in \path{backend/app/core/security.py}.
  \item \textbf{SEC-003}: Job cancellation must be scoped to the authenticated workspace before mutation.
  \item \textbf{SEC-004}: Auth endpoints use \texttt{slowapi} limits to reduce brute-force and CPU-exhaustion risk.
  \item \textbf{SEC-005}: Static Fernet key remains an accepted risk for this stage and should become KMS-backed in production.
  \item \textbf{SEC-006}: CORS is restricted to explicit origins, methods, and headers instead of a wildcard policy.
\end{itemize}

\end{document}
